\documentclass[12pt]{article}
\usepackage[T1]{fontenc}
\usepackage[utf8]{inputenc}
\usepackage[brazil]{babel}
\usepackage[margin=1in]{geometry}
\setlength{\parindent}{0pt}
\begin{document}
\section*{Tema 35}
Processo: PEDILEF 2006.63.02.012610-0/SP\\
Situacao: Julgado\\
Ramo: DIREITO PREVIDENCIÁRIO\\
Questao: Saber se em caso de contagem recíproca de tempo de contribuição, nos casos de regimes próprios de previdência dos servidores públicos, deve haver indenização de contribuições previdenciárias em caso de averbação de tempo rural.\\
Tese: A averbação de tempo de trabalho rural nos regimes próprios de previdência dos servidores públicos exige a indenização das respectivas contribuições previdenciárias, para efeitos de contagem recíproca de tempo de contribuição.\\
Fonte: https://www.cjf.jus.br/phpdoc/virtus/
\end{document}
